\section{Introduction}
\label{sec:introduction}

Every working system that stores, transmits, or processes information operates,
in the end, on what its designers call ``data.'' The term shows up in database
textbooks, programming language specifications, legal frameworks, and
information-science curricula, yet it somehow resists a precise, operational
definition. Database texts describe data as ``raw facts'' awaiting
processing~\autocite{rob2013}; information-science references describe data as
facts organized for communication~\autocite{ratzan2004}; international
guidelines call data ``facts, concepts, or instructions expressed in a
standardized way, suitable for being communicated, interpreted, or processed,
whether by people or by automated systems''~\autocite{unesco2008}; and general
dictionaries treat data as ``facts or things known, used as the basis for
inference or calculation''~\autocite{oed2023}. None of these is wrong, exactly,
but none of them is formal either. None specifies what structure a collection
of facts must have before it can be indexed, copied, compared for equality, or
serialized without ambiguity.\\
This gap is not an academic inconvenience. It shows up in production systems
as a recognizable family of bugs. Schema drift happens when two components
silently disagree about what an index or field name means, even while
operating on byte-identical data. Type confusion~\autocite{cwe843} happens when
a fixed byte sequence gets reinterpreted under a type assumption different from
the one it was originally written under. Unsound equality happens when two
values are treated as equal just because they share a representation, even
though they differ in the value set the comparison is silently assuming. And
serialization ambiguity happens when a wire format under-specifies the index
set or value set badly enough that two conformant implementations end up
disagreeing on what the data actually means. Each of these has the same root
cause: an implicit assumption about the domain, codomain, or mapping behind a
piece of data was never written down, and so it could never be checked.\\
Closing this gap takes more than just picking one informal definition over
another. It requires working through, in order, what data actually are, how
data come to mean something, and how that meaning gets encoded back into data
a machine can process. In other words, it requires a unified account of data,
representation, and interpretation. That is what this paper sets out to do.

\subsection{Contribution}

This paper makes three contributions:

\begin{enumerate}
  \item \textbf{A lexical and etymological grounding} (Section~\ref{sec:what-is-data})
        that traces the concept of data from its Latin root to contemporary
        usage in order to surface a structural implication that informal
        definitions all share but never actually state.
  \item \textbf{A formal treatment} (Section~\ref{sec:formal-definition}) of
        data as a total function from an index set to a value set, with
        explicit well-definedness conditions, a precise equality criterion,
        and a systematic account of how the usual data structures, sequences,
        tuples, multisets, fall out as special cases of this one definition.
  \item \textbf{A unified account of data, information, and representation}
        (Sections~\ref{sec:data-information-knowledge}
        and~\ref{sec:data-and-representation}) showing how the formal
        definition grounds the data--information--knowledge hierarchy, and how
        the contingency of representation follows directly once you adopt the
        function-theoretic view.
\end{enumerate}

In Section~\ref{sec:implications} we derive four engineering corollaries that
connect specific violations of the formal conditions to the failure modes
listed above, and we close in Section~\ref{sec:conclusion} with some
directions for future work.

\subsection{Scope and Assumptions}

This paper is concerned with data the way computing systems treat it: stored,
addressed, transmitted, interpreted. It is not trying to give a complete
philosophical account of the data--information--knowledge distinction, and it
does not take a side in the debate over whether data are theory-laden or
``given'' in some stronger epistemological sense. The formal treatment here is
deliberately minimal. It introduces only the mathematical structure needed to
make the engineering corollaries precise, nothing beyond that. Readers looking
for a fuller philosophical treatment should start with
Floridi~\autocite{floridi2010} and Zins~\autocite{zins2007}.
